\documentclass{beamer}
\usetheme{CambridgeUS}
% \usecolortheme{seahorse}
% \usecolortheme{whale}
\usecolortheme{spruce}
\usefonttheme{professionalfonts}

% % Change standard block colors

\setbeamerfont{block title}{series=\bfseries} % Grassetto nei titoli

% \setbeamertemplate{blocks}[default]   %%  PRESENTAZIONE SQUADRATA
\setbeamertemplate{blocks}[rounded][shadow=true]

% % 1- Block title (background and text)

\setbeamercolor{block title}{use={palette primary,palette secondary}, bg=palette secondary.bg, fg=palette primary.fg}

% % 2- Block body (background)
\setbeamercolor{block body}{use={palette primary}, bg=palette primary.bg}

% % Change alert block colors

% % 1- Block title (background and text)
\setbeamertemplate{blocks}[rounded][shadow=true]
\setbeamercolor{block title alerted}{use={palette secondary}, bg=palette secondary.fg, fg=white}

% % 2- Block body (background)
\setbeamercolor{block body alerted}{use={palette primary}, bg=palette primary.bg}


% {\usebeamercolor[fg]{palette secondary} Prova} % colore giusto per lo sfondo
% {\usebeamercolor[bg]{palette secondary} Due} % colore giusto per i titoli

\newenvironment{definizioneb}{\begin{block}{Definizione}}{\end{block}}
\newenvironment{osservazioneb}{\begin{block}{Osservazione}}{\end{block}}
\newenvironment{teoremab}{\begin{alertblock}{Teorema}}{\end{alertblock}}
\newenvironment{corollariob}{\begin{alertblock}{Corollario}}{\end{alertblock}}

% \setbeamertemplate{itemize items}[square] %%  PRESENTAZIONE SQUADRATA
\setbeamertemplate{itemize items}[circle]
\setbeamercolor{itemize item}{use={palette primary}, fg=palette primary.fg}



% PACCHETTI FONDAMENTLAI
\usepackage[T1]{fontenc}
\usepackage[utf8]{inputenc}
\usepackage[italian]{babel}
\usepackage[babel]{csquotes}
\usepackage[style=numeric, hyperref]{biblatex}
\usepackage{microtype}
\usepackage{lmodern}
\usepackage{graphicx} % inserire immagini
\usepackage{subfig} % sottofigure
\usepackage{multicol} % due colonne
% \usepackage{ulem} % sottolineare
\usepackage{lipsum} % lorem ipsum
\usepackage{xcolor} % colori in latex
\usepackage{parskip} % rimuove l'indentazione dei nuovi paragrafi %% Add parbox=false to all new tcolorbox
\usepackage{centernot}
\usepackage[outline]{contour}\contourlength{3pt}
\usepackage{fancyhdr}
\usepackage{layout}
\usepackage[intlimits]{empheq} % Riquadri colorati attorno alle equazioni
\usepackage[most]{tcolorbox} % Riquadri colorati
\usepackage{ifthen} % IFTHEN
\usepackage{rotating} %figure ruotate 

% pacchetti matematica
\usepackage[intlimits]{amsmath} 
\usepackage{amssymb}
\usepackage{amsthm}
\usepackage{yhmath}
\usepackage{dsfont}
\usepackage{mathrsfs}
\usepackage{cancel} % semplificare
\usepackage{polynom} %divisione tra polinomi
\usepackage{forest} % grafi ad albero
\usepackage{booktabs} % tabelle
\usepackage{commath} %simboli e differenziali
\usepackage{bm} %bold
\usepackage[fulladjust]{marginnote} %to use marginnote for date notes
\usepackage{arrayjobx}%array
\usepackage{mathtools}
%%%%%%%%%%%%%


%%%% CONTATORE DIMOSTRAZIONI
\newcounter{numerodidimostrazioni}
\newcommand{\mostranumerodimostrazioni}{\cleardoublepage\thispagestyle{empty}\begin{center}\null\vfill \textit{Questo corso comprende {\thenumerodidimostrazioni} dimostrazioni}\vfill\null\end{center}}
%%%%

%%%% QUIVER
\newcommand{\duepunti}{\,\mathchar\numexpr"6000+`:\relax\,}
% A TikZ style for curved arrows of a fixed height, due to AndréC.
\tikzset{curve/.style={settings={#1},to path={(\tikztostart)
    .. controls ($(\tikztostart)!\pv{pos}!(\tikztotarget)!\pv{height}!270:(\tikztotarget)$)
    and ($(\tikztostart)!1-\pv{pos}!(\tikztotarget)!\pv{height}!270:(\tikztotarget)$)
    .. (\tikztotarget)\tikztonodes}},
    settings/.code={\tikzset{quiver/.cd,#1}
        \def\pv##1{\pgfkeysvalueof{/tikz/quiver/##1}}},
    quiver/.cd,pos/.initial=0.35,height/.initial=0}

% TikZ arrowhead/tail styles.
\tikzset{tail reversed/.code={\pgfsetarrowsstart{tikzcd to}}}
\tikzset{2tail/.code={\pgfsetarrowsstart{Implies[reversed]}}}
\tikzset{2tail reversed/.code={\pgfsetarrowsstart{Implies}}}
% TikZ arrow styles.
\tikzset{no body/.style={/tikz/dash pattern=on 0 off 1mm}}
%%%%%%%%%%

%%% HEADER e FOOTER
%% HEAD and FOOT
\newcommand{\datainizion}{testo}
\newcommand{\hdrnew}[1]{
	\fancyhead{} % cancella tutti i campi
	\fancyfoot{}
	\fancyhead[LE]{#1}
	\fancyhead[RO]{\nouppercase{\leftmark}}
	% \fancyfoot[LE,RO]{\thepage}
	\fancyfoot[C]{\thepage}
	% \fancyfoot[CO,RE]{Per: Dean A. Smith}
	\renewcommand{\headrulewidth}{0.4pt}
	\renewcommand{\footrulewidth}{0pt}
	\pagestyle{fancy}
}
\newcommand{\days}[1]{\renewcommand{\datainizion}{#1}}
%%%%%%%%

%% DEFINIZIONI COMANDI MATEMATICI
\let\sin\relax %TOGLIE LA DEFINIZIONE SU "\sin"

% cambia la definizione di empty set
% ---
\let\oldemptyset\emptyset
% ---
\let\emptyset\varnothing 
% ---
% \let\emptyset\relax
% \newcommand{\emptyset}{\text{\textnormal{\O}}}
% ---

\DeclareMathOperator{\bounded}{bd}
\DeclareMathOperator{\sin}{sen}
\DeclareMathOperator{\epi}{Epi}
\DeclareMathOperator{\cl}{cl}
\DeclareMathOperator{\graph}{graph}
\DeclareMathOperator{\arcsec}{arcsec}
\DeclareMathOperator{\arccot}{arccot}
\DeclareMathOperator{\arccsc}{arccsc}
\DeclareMathOperator{\spettro}{Spettro}
\DeclareMathOperator{\nulls}{nullspace}
\DeclareMathOperator{\dom}{dom}
\DeclareMathOperator{\ar}{ar}
\DeclareMathOperator{\const}{Const}
\DeclareMathOperator{\fun}{Fun}
\DeclareMathOperator{\rel}{Rel}
\DeclareMathOperator{\altezza}{ht}
\let\det\relax %TOGLIE LA DEFINIZIONE SU "\det"
\DeclareMathOperator{\det}{det}
\DeclareMathOperator{\End}{End}
\DeclareMathOperator{\gl}{GL}
\DeclareMathOperator{\Id}{Id}
\DeclareMathOperator{\id}{Id}
\DeclareMathOperator{\I}{\mathds{1}}
\DeclareMathOperator{\II}{II}
\DeclareMathOperator{\rank}{rank}
\DeclareMathOperator{\tr}{tr}
\DeclareMathOperator{\tc}{t.c.}
\DeclareMathOperator{\T}{T}
\DeclareMathOperator{\var}{Var}
\DeclareMathOperator{\cov}{Cov}
\DeclareMathOperator{\st}{st}
\DeclareMathOperator{\mon}{Mon}
\newcommand{\trasposta}[1]{\prescript{\text{T}}{}{#1}}
\newcommand{\nonstandard}[1]{\prescript{*}{}{#1}}
\newcommand{\halo}{\mathfrak{h}}
\newcommand{\hhalo}[1]{\prescript{\halo}{}{#1}}
%%% Modo 1

% \newcommand{\sequence}[2][n \in \N]{{\langle #2\rangle}^{#1}
% }
% \newcommand{\filtersequence}[2][n \in \N]{\sequence[#1]{#2}_{\mathcal{U}}}

%%%
%%% Modo 2

\newcommand{\sequence}[2][n \in \N]{{\langle #2\rangle}_{\!#1}
}
\newcommand{\filtersequence}[2][n \in \N]{%
	{\overline{\langle #2\rangle}}_{\!#1}%
}

%%%

%%% Modo 3

% \newcommand{\sequence}[2][n \in \N]{{\langle #2\rangle}_{#1}
% }
% \newcommand{\filtersequence}[2][n \in \N]{\sequence[#1]{#2}^{\mathcal{U}}}

%%%
\newcommand{\R}{\mathds{R}}
\newcommand{\starR}{\nonstandard{\R}}
\newcommand{\borel}{\mathscr{B}}
\newcommand{\lebesgue}[1]{\mathscr{L}\left(#1\right)}
\newcommand{\media}{\mathds{E}}
\newcommand{\K}{\mathds{K}}
\newcommand{\A}{\mathds{A}}
\newcommand{\Q}{\mathds{Q}}
\newcommand{\N}{\mathds{N}}
\newcommand{\C}{\mathds{C}}
\newcommand{\Z}{\mathds{Z}}
\newcommand{\rmn}{\R^{m,n}}
\newcommand{\qo}{\hspace{1em}\text{q.o.}\,}
\newcommand{\pprec}{\mathrel{\prec\! \prec}}
\newcommand{\npprec}{\mathrel{\prec\! \nprec}}
\renewcommand{\tilde}[1]{\widetilde{#1}}
\renewcommand{\parallel}{\mathrel{/\mkern-5mu/}}
\newcommand{\cardin}[1]{\left\vert #1\right\vert}
\newcommand{\parti}[1]{\wp (#1)}
\newcommand{\diff}[1]{\operatorname{d}_{#1}}
\renewcommand{\vec}[1]{\overrightarrow{\vphantom{i}#1}}
\newcommand{\floor}[1]{\left\lfloor #1 \right\rfloor}

%tramite i prossimi due comandi posso decidere come scrivere i logaritmi naturali in tutti i documenti: ho infatti eliminato qualsiasi differenza tra "ln" e "log": se si vuole qualcosa di diverso bisogna inserire manualmente il tutto
\let\ln\relax
\DeclareMathOperator{\ln}{ln}
\let\log\relax
\DeclareMathOperator{\log}{log}
%%%%%%

%% NUOVI COMANDI
\newcommand{\straniero}[1]{\textit{#1}} %parole straniere
\newcommand{\titolo}[1]{\textsc{#1}} %titoli
\newcommand{\qedd}{\tag*{$\blacksquare$}} %qed per ambienti matemastici
\renewcommand{\qedsymbol}{$\blacksquare$} %modifica colore qed
\newcommand{\ooverline}[1]{\overline{\overline{#1}}}
\newcommand{\circoletto}[1]{\left(#1\right)^\text{o}}
%
\newcommand{\qmatrice}[1]{\begin{pmatrix}
#1_{11} & \cdots & #1_{1n}\\
\vdots & \ddots & \vdots \\
#1_{m1} & \cdots & #1_{mn}
\end{pmatrix}}
%
\newcommand{\parentesi}[2]{%
\underset{#1}{\underbrace{#2}}%
}
%
\newcommand{\norma}[1]{% Norma 
\left\lVert#1\right\rVert%
}
\newcommand{\scalare}[2]{% Scalare
\left\langle #1, #2\right\rangle
}
%%%%%

%% RESTRIZIONI
\newcommand{\referenze}[2]{
	\phantomsection{}#2\textsuperscript{\textcolor{blue}{\textbf{#1}}}
}
\def\restriction#1#2{\mathchoice
              {\setbox1\hbox{${\displaystyle #1}_{\scriptstyle #2}$}
              \restrictionaux{#1}{#2}}
              {\setbox1\hbox{${\textstyle #1}_{\scriptstyle #2}$}
              \restrictionaux{#1}{#2}}
              {\setbox1\hbox{${\scriptstyle #1}_{\scriptscriptstyle #2}$}
              \restrictionaux{#1}{#2}}
              {\setbox1\hbox{${\scriptscriptstyle #1}_{\scriptscriptstyle #2}$}
              \restrictionaux{#1}{#2}}}
\def\restrictionaux#1#2{{#1\,\smash{\vrule height .8\ht1 depth .85\dp1}}_{\,#2}} 
%%%%%%%%%%%

%% SEZIONE GRAFICA
\usepackage{tikz}
\usetikzlibrary{matrix, patterns, calc, decorations.pathreplacing, hobby, decorations.markings, decorations.pathmorphing, babel}
\usepackage{tikz-3dplot}
\usepackage{mathrsfs} %per geogebra
\usepackage{tikz-cd}
\tikzset
{
  %surface/.style={fill=black!10, shading=ball,fill opacity=0.4},
  plane/.style={black,pattern=north east lines},
  curve/.style={black,line width=0.5mm},
  dritto/.style={decoration={markings,mark=at position 0.5 with {\arrow{Stealth}}}, postaction=decorate},
  rovescio/.style={decoration={markings,mark=at position 0.5 with {\arrow{Stealth[reversed]}}}, postaction=decorate}
}
\usepackage{pgfplots} % stampare le funzioni
	\pgfplotsset{/pgf/number format/use comma,compat=1.15}
	%\pgfplotsset{compat=1.15} %per geogebra
	\usepgfplotslibrary{fillbetween, polar}
%%%%%%

%% CITAZIONI
\usepackage{lineno}

\newcommand{\citazione}[1]{%
  \begin{quotation}
  \begin{linenumbers}
  \modulolinenumbers[5]
  \begingroup
  \setlength{\parindent}{0cm}
  \noindent #1
  \endgroup
  \end{linenumbers}
  \end{quotation}\setcounter{linenumber}{1}
  }
%%%%%%


\usepackage{hyperref}
\hypersetup{%
	pdfauthor={Davide Peccioli},
	pdfsubject={Appunti UniTO},
	allcolors=black,
	citecolor=black,
	colorlinks=true, 
	bookmarksopen=true}


\title[Davide Peccioli]{Flusso di campi vettoriali su varietà e parentesi di Lie: un approccio non standard.}
\author[]{Candidato: Davide Peccioli\\ Relatore: prof. Luca Motto Ros}
\date{11 novembre 2024}

\begin{document}


\begin{frame}
    \maketitle
\end{frame}


\begin{frame}
    \frametitle{Prerequisiti}

    Definizione di limite.

    Elementi interni di $\nonstandard{V}(X) $.

\end{frame}

\begin{frame}
    % \frametitle{}

    \begin{definizioneb}
        Per $ r, s \in \starR $, si adotterà la seguente notazione: \begin{itemize}
            \item $ r \prec s $ se $ r = a \cdot s $ per qualche $ a $ \emph{finito};
            \item $ r \pprec s $ se $ r = \lambda \cdot s $ per qualche $ \lambda $ infinitesimo.
        \end{itemize}
    \end{definizioneb}

    Questa definizione si estende naturalmente ad elementi di $ \nonstandard{(\R^n)} $. In particolare, se $ v \in \nonstandard{(\R^n)} $ e $ s \in \starR $, si adotta la seguente notazione: \begin{itemize}
        \item $ v \prec s $ se $ \nonstandard{\norm{v}} \prec s $ per qualche norma $ \norm{ \cdot } $ su $ \R^{n} $;
        \item $ v \pprec s $ se $ \nonstandard{\norm{v}} \pprec s $ per qualche norma $ \norm{ \cdot } $ su $ \R^{n} $;
    \end{itemize}

    Se $ v, w \in \nonstandard{(\R^n)} $, si scriverà $ v\approx w $ se $ v-w \pprec 1 $.

\end{frame}

\begin{frame}
    % \frametitle{<title>}

    Sia $ M $ una varietà differenziabile fissata, e sia $\lambda \in \starR$ un infinitesimo fissato. 

    \begin{definizioneb}
        Per $ p \in M $, si definisce $ \halo(p) \subseteq \nonstandard{M} $ l'insieme dei punti $ x \in \nonstandard{M} $ per i quali vi è un intorno coordinato $ U $ di $ p $ tale che $ x \in \nonstandard{U} $, e $ x \approx p $ in quelle date coordinate.
    \end{definizioneb}

    Per ogni $ A \subseteq M $ si definisce \emph{l'alone} di $ A $ come \begin{equation*}
        \hhalo{A} = \bigcup_{a \in A}\halo(A). 
    \end{equation*}I punti di $ \hhalo{M} $ si dicono \emph{quasi standard} e, per ogni $ a \in \hhalo{M} $, si definisce la \emph{parte standard} di $ a $, denotata con $ \st(a) $, come l'unico $ p \in M $ tale per cui $ a \in \halo(p) $. 

    % Poiché $ M $ è compatta, allora $ \hhalo{M} = \nonstandard{M} $. % Notare il discorso della compattezza
\end{frame}

\begin{frame}
    % \frametitle{<title>}

    \begin{definizioneb}
        Sia $ a \in \hhalo{M} $. Un \emph{prevettore} in $ a $ è una coppia $ (a,x) $ con $ x \in \hhalo{M} $ tale che per ogni funzione differenziabile $ f: M \longrightarrow \R $, \[
        \nonstandard{f}(x)-\nonstandard{f}(a)\prec \lambda
    \]Si indicherà con $ P_{a}M  $ l'insieme dei prevettori in $ a $.
    \end{definizioneb}

    Si definisce una relazione di equivalenza $ \equiv $ in $ P_{a}M  $ come segue: \[
        (a,x) \equiv (a,y) \qquad \text{se e solo se}\qquad \nonstandard{f}(y)-\nonstandard{f}(x)\pprec \lambda
    \]per ogni funzione differenziabile $ f: M \longrightarrow \R $.

    L'insieme delle classi di equivalenza $ P_{a}M/\equiv  $ si denoterà con $ T_{a}M  $, e la classi di equivalenza di $ (a,x) $ con $ \vec{ax} $. Inoltre, $ T_{a}M $ è uno spazio vettoriale reale $ n $-dimensionale, indipendente dalla scelta delle coordinate.

\end{frame}

\begin{frame}
    % \frametitle{<title>}

        Sia $ a \in \hhalo{M} $. e sia $ f:M \longrightarrow \R $ una funzione differenziabile. \begin{itemize}
            \item Un prevettore $ (a,x) \in P_{a}M  $ agisce su una funzione differenziabile $ f: M \longrightarrow \R $ come segue: \[
                (a,x)\,f \coloneqq \frac{1}{\lambda}\left(\nonstandard{f}(x)-\nonstandard{f}(a)\right) \in \starR
            \]che è finito per definizione di prevettore.
            \item Un vettore $ \vec{ax} \in T_{a}M  $ agisce come derivazione su $ f $ come segue: \[
                \vec{ax} \,f = \st\left((a,x)\,f\right) \in \R.
            \] 
        \end{itemize}
    
    \begin{teoremab}
        Ogni vettore classico $ X $ tangente ad $ M $ in $ a \in M$ definisce un unico vettore $ \vec{ax} \in T_{a} M  $ dove, in coordinate, $ x = a + \lambda X $.
    \end{teoremab}

\end{frame}

\begin{frame}
    \frametitle{Campi di prevettori}
    \begin{definizioneb}
        Un \emph{campo di prevettori} su $ \nonstandard{M} $ è una mappa \emph{interna} $ F: \nonstandard{M} \longrightarrow \nonstandard{M} $ tale che $ \left(a, F(a)\right) \in P_{a}M  $ per ogni $ a \in \hhalo{M} $. 
        % Equivalentemente, \begin{equation*}
        %     F(a)-a\prec \lambda\qquad \forall\,a \in \hhalo{M}.
        % \end{equation*}
    \end{definizioneb}

    Se $ G: \nonstandard{M} \longrightarrow \nonstandard{M} $ è un altro campo di prevettori, $ F $ e $ G $ si dicono \emph{equivalenti} (e si scrive $ F\equiv G $) se \begin{equation*}
        \left(a,F(a)\right) \equiv \left(a,G(a)\right)\qquad \forall\,a \in \hhalo{M}
    \end{equation*}

    \begin{definizioneb}
        Un campo di prevettori \emph{locale} è una mappa interna $F: \nonstandard{U} \longrightarrow \nonstandard{V}$ tale che  $ \left(a, F(a)\right) \in P_{a}M  $ per ogni $ a \in \hhalo{U} $, con $ U \subseteq V \subseteq M $ aperti.
    \end{definizioneb}

\end{frame}

\begin{frame}
    \frametitle{Relazione con i campi vettoriali}

    Inoltre, $ F $ \emph{realizza} un campo vettoriale classico $ X $ se per ogni funzione differenziabile $ h : U \longrightarrow \R $ si ha che \begin{equation*}
        X\,h(a) = \vec{a\, F(a)}\,h,\qquad \forall\, a \in U.
    \end{equation*} 

    \begin{teoremab}
        Ogni campo vettoriale $ X $ su una varietà ammette localmente un campo locale di prevettori che realizza $ X $.
    \end{teoremab}

    Il campo di prevettori definito in coordinate come \begin{equation*}
        F(a) = a + \lambda \cdot X(a) \forall\, a \in \nonstandard{M}
    \end{equation*}realizza $ X $.     

\end{frame}

\begin{frame}
    % \frametitle{<title>}

    \begin{definizioneb}
        Un campo di prevettori $ F $ è di classe $ D^{1} $ se per ogni $ a \in \hhalo{M} $ e $ v \in \nonstandard{(\R^{n})}$ con $ v \prec \lambda $ si ha, in coordinate \begin{equation*}
            F(a)-F(a+v)+v \prec \lambda\norma{v}.
        \end{equation*}
    \end{definizioneb}

    Si dimostra che $ F $ è invertibile, e $ F^{-1} $ è di classe $ D^{1} $.

    % Questa definizione è motivata dal modo in cui la condizione di Lipschitzianità su un campo vettoriale si trasmette sul campo di prevettori che lo realizza. Se $ X $ è un campo vettoriale Lipschitziano, allora il campo di prevettori $ F $ tale che in coordinate $ F(a) = a + \lambda \cdot X $ è di classe $ D^{1} $.

    Similmente, un campo di prevettori $ G $ si dice di classe $ D^{2} $ se, per ogni $ a \in \hhalo{M} $ e per ogni $ v,w \in \nonstandard{(\R^n)} $, con $ v,w \prec \lambda  $, si ha, in coordinate, \begin{equation*}
        G(a)-G(a+v)-G(a+w)+G(a+v+w) \prec \lambda\,\norma{v}\norma{w }. 
    \end{equation*}Inoltre, $ G $ è di classe $ D^{1} $.
\end{frame}

\begin{frame}
    % \frametitle{<title>}

    \begin{definizioneb}
        Sia $ F: \nonstandard{M} \longrightarrow \nonstandard{M} $ un campo di prevettori, e, per $ 0\le t \in \starR $, sia $ n = n(t) = \floor{t/\lambda} $. Per ogni $ a \in \nonstandard{M} $ si definisce il \emph{flusso di $ F $ al tempo $ t $}, $ F_{t}(a) = F^{n}(a) $. 
    \end{definizioneb}

    Sia $ \operatorname{Map}(M) $ l'insieme di tutte le mappe $ f: M \longrightarrow M $ e sia $ \Phi $ la mappa\begin{align*}
    \Phi:\operatorname{Map}(M) \times \N &\longrightarrow \operatorname{Map}(M) \\
    (f,m) &\longmapsto \underset{m\text{ volte}}{\underbrace{f\circ \dots\circ f}}.
    \end{align*}Presa la sua estensione non standard $ \nonstandard{\Phi}: \nonstandard{\operatorname{Map}(M)}\times \nonstandard{\N} \longrightarrow \nonstandard{\operatorname{Map}(M)} $, si definisce $ F^{n} \coloneqq \nonstandard{\Phi}(F,n) $.

    % Se $ F $ è invertibile (e questo accade per campi di prevettori di classe $ D^{1} $), si definisce $ F_{t}$ per $ t\le 0 $ come $ (F^{-1})_{-t} $.
\end{frame}

\begin{frame}
    \frametitle{Flusso standard}

    Sia $ F $ un campo di prevettori di classe $ D^{1} $, invertibile. % Questa ipotesi non è restrittiva, poiché è possibile ampliare la teoria dei campi di prevettori localmente, e tutti i campi di prevettori locali sono localmente invertibili.
    Si dimostra che esiste esiste $ U \subseteq M $ aperto ed esiste un intervallo $ [-T,T] \subseteq \R $ tale per cui la seguente definizione è ben posta.
    
    \begin{definizioneb}
        Sia $ t \in [-T,T] $. Il \emph{flusso standard} indotto da $ F $ al tempo $ t $ è l'applicazione \begin{align*}
            h_{t}^{F}: U &\longrightarrow M \\ 
            x & \mapsto \st\left(F_{t}(x) \right).
        \end{align*}
    \end{definizioneb}

    Inoltre, se $ X $ è un campo vettoriale classico e $ F $ è il campo di prevettori che realizza $ X $, $ F(a) = a + \lambda\,X $, allora $ h_{t}^{F}$ è \emph{esattamente} il flusso di $ X $ nel senso classico. 

\end{frame}

\begin{frame}
    \frametitle{Parentesi di Lie}

    \begin{definizioneb}
        Siano $ F,G $ campi di prevettori di classe $ D^{1} $. La parentesi di Lie di $ F $ e $ G $ è il campo di prevettori \begin{equation*}
            [F,G] = \left(G^{-1}\circ F^{-1}\circ G\circ F\right)^{\floor{1/\lambda}}.
        \end{equation*}
    \end{definizioneb}

    Se inoltre $ F $ e $ G $ sono di classe $ D^{2} $, allora, per ogni $ a \in \nonstandard{M} $ si ha \begin{equation*}
        [F,G](a) \equiv a + \frac{1}{\lambda} \left(G^{-1}\circ F^{-1}\circ G\circ F(a)-a\right)
    \end{equation*}

    % \begin{teoremab}
    %     Siano $ X,Y $ due campi vettoriali di classe $ C^{2} $, e sia $ [X,Y]_{\text{cl}} $ la loro parentesi di Lie classica. 

    %     Siano $ F,G $ due campi di prevettori che realizzino $ X,Y $ rispettivamente. Allora $ [F,G] $ realizza $ [X,Y]_{\text{cl}} $.
    % \end{teoremab}

\end{frame}

\begin{frame}
    \begin{teoremab}
        Siano $ X,Y $ due campi vettoriali di classe $ C^{2} $, e sia $ [X,Y]_{\text{cl}} $ la loro parentesi di Lie classica. 

        Siano $ F,G $ due campi di prevettori che realizzino $ X,Y $ rispettivamente. Allora $ [F,G] $ realizza $ [X,Y]_{\text{cl}} $.
    \end{teoremab}

    Ricordando che i flussi $ h_{t}^{F} $ e $ h_{t}^{G} $ coincidono con i flussi di $ X $ e $ Y $ nel senso classico, e che, per la parentesi di Lie vale il ben noto risultato \begin{equation*}
        [X,Y]_{\text{cl}} (p) = \lim_{t\to 0} \frac{1}{t^{2}}\left(\left(h_{t}^{G} \right)^{-1} \circ \left(h_{t}^{F} \right)^{-1}\circ h_{t}^{G}\circ h_{t}^{F}(p) - p \right).
    \end{equation*}Sfruttando la definizione di limite nell'analisi non standard, si ha \begin{equation*}
        [X,Y]_{\text{cl}} (p) \approx \frac{1}{\lambda^{2}}\left(\left(h_{\lambda}^{G} \right)^{-1} \circ \left(h_{\lambda}^{F} \right)^{-1}\circ h_{\lambda}^{G}\circ h_{\lambda}^{F}(p) - p \right).
    \end{equation*}

\end{frame}

\begin{frame}

    \begin{teoremab}
        Siano $ X,Y $ due campi vettoriali di classe $ C^{2} $, e sia $ [X,Y]_{\text{cl}} $ la loro parentesi di Lie classica. 

        Siano $ F,G $ due campi di prevettori che realizzino $ X,Y $ rispettivamente. Allora $ [F,G] $ realizza $ [X,Y]_{\text{cl}} $.
    \end{teoremab}
    
    Dunque $ [X,Y]_{\text{cl}} $ è realizzato dal campo di prevettori \begin{equation*}
        p\mapsto p + \frac{1}{\lambda}\left(\left(h_{\lambda}^{G} \right)^{-1} \circ \left(h_{\lambda}^{F} \right)^{-1}\circ h_{\lambda}^{G}\circ h_{\lambda}^{F} (p)- p \right)
    \end{equation*}che è equivalente al campo $ [F,G] $.
    
    Quindi $ [F,G] $ realizza $ [X,Y]_{\text{cl}} $.
\end{frame}

\begin{frame}
    \frametitle{Conclusioni}

    \begin{teoremab}
        Siano $ F,G $ due campi di prevettori di classe $ D^{2} $. 

        Allora  $ [F,G]\equiv I $, dove $ I $ è l'identità, se e solo se esiste $ 0<T \in \R $ tale che per ogni $ 0\le t,s\le T $:\[h_{t}^{F}\circ h_{s}^{G} = h_{s}^{G}\circ h_{t}^{F}.\]
    \end{teoremab}
    \begin{corollariob}
        Siano $ X,Y $ dei campi vettoriali classici di classe $ C^{2} $. 
    
        Allora i flussi di $ X $ e di $ Y $ commutano se e solo se $ [X,Y]_{\text{cl}}=0 $.
    \end{corollariob}

    Il corollario segue banalmente dal teorema poiché il campo di prevettori identitario realizza il campo vettoriale nullo, e la parentesi di Lie realizza la parentesi di Lie standard. 

\end{frame}

\begin{frame}
    
La dimostrazione del teorema verte sulla natura combinatorica di $ h^{F}_t $ e $ h_{s}^{G}  $. Sia $ n \coloneqq\floor{t/\lambda} $ e $ m\coloneqq \floor{s/\lambda } $. Se \begin{equation*}
    \forall\, a \in \nonstandard{M}:\quad F^{n}\circ G^{m} (a) \approx G^{m}\circ F^{n}(a).
\end{equation*}allora $ h^{F}_t\circ h^{G}_s = h^{G}_s\circ h^{F}_s$ applicando ad entrambi i membri la funzione $ \st $. 

\end{frame}

\end{document}